\section*{Capítulo \ref{ch:b2:measurement-statistics} --- Estatística da medida química}

\begin{solution}{exo:b2:measurement-statistics:1}
$\bar V = 62.41/5 = \qty{12.482}{mL}$; $\sum(V_i - \bar V)^2 = 0.00328$, $s = \sqrt{0.00328/4} \approx
\qty{0.029}{mL}$; $s/\sqrt5 \approx \qty{0.013}{mL}$.
\end{solution}

\begin{solution}{exo:b2:measurement-statistics:2}
$t = 2.776$ para 4 graus de liberdade: $12.482 \pm 2.776 \times 0.0128$, isto é, $(12.48 \pm 0.04)$~mL.
\end{solution}

\begin{solution}{exo:b2:measurement-statistics:3}
As incertezas relativas da massa e do volume são 0.04~\% e 0.06~\%; combinadas em
quadratura,
\[\frac{u(c)}{c} = \sqrt{\Bigl(\frac{0.0002}{0.5000}\Bigr)^2 + \Bigl(\frac{0.15}{250.0}\Bigr)^2}
\approx 0.07~\%.\]
\end{solution}

\begin{solution}{exo:b2:measurement-statistics:4}
$\bar x = 1.5$, $\bar y = 0.31$, $S_{xx} = 5$, $S_{xy} = 0.99$: $b = 0.198$, $a = 0.31 - 0.198 \times 1.5
= 0.013$.
\end{solution}

\begin{solution}{exo:b2:measurement-statistics:5}
$|10.12 - 10.00|/(0.08/\sqrt5) = 0.12/0.036 \approx 3.4 > 2.776$: a diferença é significativa, o
método é enviesado (de cerca de \qty{+0.12}{mg/L}).
\end{solution}

\begin{solution}{exo:b2:measurement-statistics:6}
$\mathrm{pH} = -\log_{10}c$, $\mathrm d\,\mathrm{pH}/\mathrm dc = -1/(c\ln10)$:
$u(\mathrm{pH}) = (u(c)/c)/\ln10 = 0.02/2.303 \approx 0.009$.
\end{solution}

\begin{solution}{exo:b2:measurement-statistics:7}
$x_{\mathrm{LD}} = 3 \times 0.0012/0.0098 \approx \qty{0.37}{\micro g/L}$; $x_{\mathrm{LQ}} = 10
\times 0.0012/0.0098 \approx \qty{1.2}{\micro g/L}$.
\end{solution}

\begin{solution}{exo:b2:measurement-statistics:8}
A reta dá \qty{3.11}{mg/L} nas soluções medidas; a amostra tinha sido diluída cinco vezes:
$3.11 \times 50.0/10.0 \approx \qty{15.6}{mg/L}$.
\end{solution}

\begin{solution}{exo:b2:measurement-statistics:9}
\qty{0.05}{mg/L} é a \omterm{def:b2:measurement-statistics:validation}{repetibilidade}, \qty{0.15}{mg/L} a \omterm{def:b2:measurement-statistics:validation}{reprodutibilidade}. Cada laboratório tem seu
próprio pequeno \omterm{def:b2:measurement-statistics:validation}{viés} (calibração, padrões, instrumento); esses vieses diferem entre laboratórios e se
somam à dispersão, de modo que a \omterm{def:b2:measurement-statistics:validation}{reprodutibilidade} é sempre a maior.
\end{solution}

\begin{solution}{exo:b2:measurement-statistics:10}
$\partial S/\partial a = 0$ se escreve $\sum(y_i - a - bx_i) = 0$: a soma dos \omterm{def:b2:measurement-statistics:calibration}{resíduos} é nula. Dividindo por
$n$: $\bar y = a + b\bar x$, logo $(\bar x, \bar y)$ está na reta.
\end{solution}

\begin{solution}{exo:b2:measurement-statistics:11}
Uma etapa:
\[\sqrt{(0.006/1.000)^2 + (0.08/100.0)^2} \approx 0.61~\%.\]
Uma etapa de dez vezes:
\[\sqrt{(0.02/10.00)^2 + (0.08/100.0)^2} \approx 0.22~\%,\]
e duas delas em quadratura dão $0.22\sqrt2 \approx 0.30~\%$. A via em duas etapas é duas vezes mais precisa: a pequena pipeta domina a via em uma
etapa.
\end{solution}

\begin{solution}{exo:b2:measurement-statistics:12}
$z = 1.6/\sqrt{0.4^2 + 0.5^2} \approx 2.5 > 2$: não compatíveis; um laboratório (pelo menos) tem um \omterm{def:b2:measurement-statistics:validation}{viés}.
Com as incertezas expandidas ($k = 2$), o primeiro dá $9.6 \pm 0.8$, que inclui 10; o segundo
$11.2 \pm 1.0$, inteiramente acima de 10. Antes de qualquer veredicto sobre a água, o desacordo deve ser
resolvido, por exemplo com um material de referência.
% ledger: who:lead
\end{solution}

\begin{solution}{pb:b2:measurement-statistics:1}
\textbf{1.} O ácido muda a atomização e a resposta; padrões e amostras devem ter a mesma
matriz para que a inclinação se aplique.
\textbf{2.} $\bar x = 10$, $\bar y = 0.1004$, $S_{xx} = 250$, $S_{xy} = 2.445$: $b = 0.00978$ por
\unit{\micro g/L}, $a = 0.1004 - 0.0978 = 0.0026$.
\textbf{3.} $+0.0004$, $-0.0005$, $+0.0006$, $-0.0013$, $+0.0008$.
\textbf{4.} Os sinais alternam, nenhuma curvatura: a reta é adequada em 0--\qty{20}{\micro g/L}.
\textbf{5.} $s_{y/x} = \sqrt{3.1 \times 10^{-6}/3} \approx 0.00102$.
\textbf{6.} O branco (ácido, água, forno) dá ele próprio uma pequena absorbância.
\textbf{7.} A sensibilidade: absorbância por \unit{\micro g/L} de chumbo.
\textbf{8.} $\bar y_0 = 0.1010$; $s = 0.0030$.
\textbf{9.} $x_0 = (0.1010 - 0.0026)/0.00978 \approx \qty{10.06}{\micro g/L}$.
\textbf{10.} $s_{x_0} = (0.00102/0.00978)\sqrt{1/3 + 1/5 + (0.1010 - 0.1004)^2/(0.00978^2 \times 250)}
\approx 0.104 \times 0.730 \approx \qty{0.076}{\micro g/L}$.
\textbf{11.} $n - 2 = 3$; $t = 3.182$.
\textbf{12.} $10.06 \pm 3.182 \times 0.076 = (10.06 \pm 0.24)$~\unit{\micro g/L}.
\textbf{13.} O termo $1/m$ cai de 1 para $1/3$: $s_{x_0}$ diminui cerca de um terço.
\textbf{14.} $f = 50.50/50.0 = 1.010$; $10.06 \times 1.010 \approx \qty{10.16}{\micro g/L}$.
\textbf{15.} $f = 1 + V_2/V_1$: $u^2(f) = (u(V_2)/V_1)^2 + (V_2u(V_1)/V_1^2)^2 = (0.005/50.0)^2 +
(0.50 \times 0.03/2500)^2$, $u(f) \approx 1.0 \times 10^{-4}$, ou 0.01~\% em valor relativo.
\textbf{16.} Não: 0.01~\% contra $0.076/10.06 \approx 0.75~\%$.
\textbf{17.} $s_{\mathrm{branco}} \approx 0.00115$; $x_{\mathrm{LD}} = 3 \times 0.00115/0.00978 \approx
\qty{0.35}{\micro g/L}$.
\textbf{18.} $10 \times 0.00115/0.00978 \approx \qty{1.2}{\micro g/L}$.
\textbf{19.} Sim, de longe.
\textbf{20.} Sim: de 9.92 a \qty{10.40}{\micro g/L} na água de torneira.
\textbf{21.} $(0.104 - 0.0026)/0.00978 \times 1.010 \approx \qty{10.5}{\micro g/L}$ e, lida sozinha,
um ``acima do valor-guia'' afirmado com segurança que os dados não sustentam.
\textbf{22.} A regra usada para construir o intervalo captura o valor verdadeiro em 95~\% das séries
às quais é aplicada; não é uma probabilidade sobre este valor em particular.
\textbf{23.} Mais leituras em réplica e mais padrões perto de \qty{10}{\micro g/L} (os termos $1/m$ e
$1/n$), ou um método mais sensível; o intervalo só diminui lentamente, como $1/\sqrt{\text{número}}$.
\textbf{24.} Analisando uma água de referência certificada, ou dopando a amostra com uma quantidade
conhecida de chumbo e verificando a recuperação.
\textbf{25.} Perto de \qty{10}{\micro g/L}, a incerteza dos métodos de rotina é de alguns por cento, como
aqui: um valor-guia mais baixo não poderia ser verificado de modo confiável por laboratórios comuns.
\textbf{26.} $\boldsymbol{(10.16 \pm 0.24)}$~\unit{\micro g/L} a 95~\%: o intervalo contém
\qty{10}{\micro g/L}, de modo que a medida não mostra nem que a água excede o valor-guia nem que o
respeita com folga.
\end{solution}
